\subsection{Irreducible Quadratic Factors} 
There are two ways we can discover an irreducible quadratic factor of a polynomial:
\begin{enumerate*}
\item We may factor out enough linear factors that the remaining factor is quadratic, and discover that it is irreducible.
\item We may \makeblank[1in] or be \makeblank[1in] an irreducible quadratic factor. (In real life, this would typically be based on how we got the polynomial.)
\end{enumerate*}
Suppose that we've discovered an irreducible quadratic factor. How can we use it?
\begin{example}
The polynomial $x^4-4x^3+3x^2+2x-6$ has $x^2-2x+2$ as an irreducible quadratic factor. Use this fact to factor it fully.
\begin{lpic}{}{12}
\regtextleft{0.5}{-1}{First, we divide by};
\regtextleft{0.5}{-10}{We can factor};
\regtextleft{10}{-10}{using basic techniques.};
\regtextleft{0.5}{-12}{$x^4-4x^3+3x^2+2x-6={}$};
\end{lpic}
\end{example}